\chapter{Grundlagen}
\label{ch:02_theoretical-foundations}

Dieses Kapitel dient der begrifflichen Fundierung der Arbeit und klärt zentrale Konzepte, die für das Verständnis nutzerzentrierter Gestaltung interner Wissens- und Kollaborationsplattformen erforderlich sind. Ziel ist es, ein einheitliches Begriffsverständnis zu schaffen und konzeptionelle Abgrenzungen vorzunehmen, auf denen die nachfolgenden Analysen aufbauen \cite{500:ISO9241-210,303:Rosenfeld2015,408:Hassenzahl2010,406:Albers2015}.
\vspace{1em}

\section{Begriffsklärungen}
\label{sec:02-01_definitions-and-concepts}
\section{Relevante UX-Methoden für interne Plattformen}
\label{sec:02-02_ux-methods}

\subsection{Personas und User Journeys}
\label{subsec:02-02-01_personas-and-journeys}
\subsection{Design Patterns und Pattern Libraries}
\label{subsec:02-02-02_design-patterns}
\subsection{UX Writing und Microinteractions}
\label{subsec:02-02-03_ux-writing-and-microinteractions}
\section{Informationsarchitektur als kognitive Orientierungshilfe}
\label{sec:02-03_information-architecture}

\Gls{ia}  bildet das strukturelle Fundament digitaler Systeme und übernimmt insbesondere in komplexen Informationsumgebungen eine zentrale Orientierungsfunktion \cite{402:Rosenfeld2015}. Sie bestimmt, wie Inhalte gruppiert, benannt und miteinander verknüpft sind und beeinflusst damit maßgeblich, wie Nutzende Informationen wahrnehmen, interpretieren und nutzen \cite{402:Rosenfeld2015,403:Westerman2014}. In internen Wissens- und Kollaborationsplattformen fungiert \Gls{ia} als kognitive Landkarte, die es ermöglicht, umfangreiche Informationsbestände effizient zu erschließen und in bestehende Arbeitsprozesse zu integrieren \cite{402:Rosenfeld2015,300:FraunhoferIESE2018}.

Eine nutzerzentrierte \Gls{ia} orientiert sich nicht primär an organisatorischen Zuständigkeiten oder technischen Systemgrenzen, sondern an den Aufgaben, Zielen und Denkmodellen der Nutzenden \cite{402:Rosenfeld2015,905:Norman1983}. Indem sie mentale Modelle, Informationssuchstrategien und begrenzte kognitive Ressourcen berücksichtigt, leistet sie einen wesentlichen Beitrag zur Reduktion von Komplexität und unterstützt eine nachhaltige Nutzung digitaler Plattformen \cite{105:Sweller1988,414:Johnson2021}. Praxisorientierte Leitlinien wie das Prinzip „Don’t make me think“ unterstreichen dabei, dass verständliche Struktur, eindeutige Benennungen und klare Navigationspfade zentrale Voraussetzungen für schnelle Orientierung und fehlerarme Interaktion sind \cite{415:Krug2014,414:Johnson2021}.

\subsection{Navigationsverhalten und mentale Modelle}
\label{subsec:02-03-01_navigation-and-mental-models}
\subsection{Reduktion kognitiver Last durch Struktur}
\label{subsec:02-03-02_cognitive-load}

\Gls{cognitiveload} entsteht, wenn Nutzende Informationen aktiv verarbeiten, Entscheidungen treffen oder Zusammenhänge erschließen müssen und dabei die begrenzte Kapazität des Arbeitsgedächtnisses beansprucht wird \cite{105:Sweller1988,414:Johnson2021}. Eine ungünstig gestaltete \Gls{ia} kann diese Belastung erhöhen, etwa durch tief verschachtelte Navigationsstrukturen, uneinheitliche Benennungen oder redundante Inhalte, die extrinsische kognitive Belastung erzeugen. Darunter wird jene zusätzliche mentale Anstrengung verstanden, die nicht aus der eigentlichen Arbeitsaufgabe entsteht, sondern durch die Struktur und Darstellung des Systems verursacht wird und damit Ressourcen von den eigentlichen Arbeitsaufgaben abzieht \cite{105:Sweller1988,414:Johnson2021}.

Eine strukturierte und klar gegliederte \Gls{ia} reduziert kognitive Last, indem sie Orientierung bietet, Auswahlmöglichkeiten begrenzt und relevante Inhalte sichtbar priorisiert \cite{414:Johnson2021,415:Krug2014}. Prinzipien wie sinnvolle Gruppierung, Hierarchisierung und progressive Offenlegung unterstützen Nutzende dabei, Informationen schrittweise zu erfassen und Aufgaben effizient zu bearbeiten, ohne von irrelevanten Details abgelenkt zu werden \cite{415:Krug2014}. Dadurch verbessert sich nicht nur die Nutzbarkeit, sondern auch die wahrgenommene Nützlichkeit des Systems, da Aufgaben mit geringerem kognitivem Aufwand und höherer Erfolgswahrscheinlichkeit erledigt werden können \cite{102:Davis1989,300:FraunhoferIESE2018}.

\subsection{Gestaltungsprinzipien der Informationsarchitektur}
\label{subsec:02-03-03_ia-design-principles}
\section{UX und digitale Transformation}
\label{sec:02-04_ux-in-digital-transformation}

\subsection{UX als Treiber organisationaler Veränderung}
\label{subsec:02-04-01_ux-as-change-driver}
\subsection{Transformationsdesign und Change Communication}
\label{subsec:02-04-02_change-communication}